\section{Đánh giá quá trình thực hiện đồ án}

Trong quá trình thực hiện đồ án, nhóm đã đạt được các kết quả so với mục tiêu đề ra:

\subsection{Các kết quả đạt được}

\begin{itemize}

    \item \textbf{Làm rõ các kiến thức nền tảng:}
    \begin{itemize}
        \item \textbf{Kiến thức về Recommender Systems:} Hiểu được nguyên lý, các thành phần, cơ chế hoạt động của các phương pháp gợi ý truyền thống (Content-based, Collaborative Filtering) và hạn chế của chúng.
        \item \textbf{Session-based Recommendation:} Nắm vững đặc điểm của bài toán gợi ý theo phiên, các thách thức kỹ thuật và yêu cầu xử lý thời gian thực.
        \item \textbf{Container Orchestration:} Hiểu sâu về Kubernetes, các đối tượng trừu tượng hóa (Pod, Deployment, Service, Ingress) và cơ chế điều phối container.
    \end{itemize}

    \item \textbf{Phân tích và đánh giá hệ thống Florus cũ:}
    \begin{itemize}
        \item Đánh giá chi tiết hiện trạng vận hành của hệ thống Florus trên nền tảng Docker Compose.
        \item Xác định các hạn chế cốt lõi: thiếu khả năng điều phối, Service Discovery thủ công, không có High Availability cho microservices, CI/CD đứt gãy.
        \item Phân tích sự xung đột kiến trúc giữa HDFS và Kubernetes, đề xuất MinIO thay thế.
    \end{itemize}

    \item \textbf{Đề xuất và thiết kế giải pháp Kubernetes:}
    \begin{itemize}
        \item Đề xuất lộ trình chuyển đổi 2 giai đoạn: Thiết lập nền tảng Cluster và Quy trình triển khai ứng dụng.
        \item Xây dựng chiến lược tái cấu trúc từ mô hình Imperative sang Declarative.
        \item Thiết kế bảng ánh xạ chuyển đổi tài nguyên từ Docker Compose sang Kubernetes Manifests.
        \item Xây dựng mô hình Umbrella Chart với phân tầng (Infrastructure, Platform, Monitoring) để giải quyết vấn đề quản lý phức tạp.
    \end{itemize}

    \item \textbf{Thiết kế kiến trúc hệ thống E-commerce tích hợp Florus:}
    \begin{itemize}
        \item Đề xuất kiến trúc Microservices với sự tách biệt giữa tầng E-commerce và tầng xử lý dữ liệu Florus.
        \item Thiết kế các luồng thực thi chính: Luồng nghiệp vụ E-commerce và Luồng dữ liệu (Data Pipeline).
        \item Thiết kế cơ chế Hybrid để giải quyết Cold-start (AI Model + Popular Engine).
        \item Xây dựng mô hình dữ liệu Polyglot Persistence với MongoDB, Neo4j, Redis và MinIO.
    \end{itemize}

    \item \textbf{Thực nghiệm và chứng minh tính khả thi:}
    \begin{itemize}
        \item Thiết lập thành công môi trường Minikube với đầy đủ các addons cần thiết.
        \item Triển khai thành công một số các thành phần Infrastructure (MongoDB, MinIO).
        \item Triển khai thành công 6 microservices của Florus và kiểm tra connectivity.
        \item Xây dựng Umbrella Helm Charts hoạt động đúng với mô hình phân tầng.
    \end{itemize}
\end{itemize}

\subsection{Các hạn chế và thách thức}

Tuy nhiên, nhóm cũng gặp một số hạn chế trong quá trình hiện thực:

\begin{itemize}
    \item \textbf{Về mô hình AI:}
    \begin{itemize}
        \item Chưa huấn luyện mô hình Session-based Recommendation trên dữ liệu thực.
        \item Chưa tích hợp mô hình AI vào Serving Service để thử nghiệm end-to-end.
    \end{itemize}

    \item \textbf{Về hạ tầng:}
    \begin{itemize}
        \item Thử nghiệm mới chỉ thực hiện trên Minikube (single-node), chưa triển khai trên multi-node cluster.
        \item Chỉ mới hiện thực được một phần của hệ thống Florus, chưa chuyển đổi hoàn thiện hệ thống lên nền tảng K8s.
        \item Chưa cấu hình High Availability cho các database (chỉ chạy single instance).
        \item Chưa thiết lập pipeline CI/CD tự động với GitOps.
    \end{itemize}

    \item \textbf{Về đánh giá hiệu năng:}
    \begin{itemize}
        \item Chưa thực hiện benchmark so sánh hiệu năng giữa kiến trúc cũ và mới.
        \item Chưa có dữ liệu thực tế để đánh giá throughput và latency của hệ thống.
    \end{itemize}
\end{itemize}

\subsection{Bảng đối chiếu mục tiêu}

\begin{table}[H]
\centering
\caption{Đối chiếu kết quả với mục tiêu đề ra}
\label{tab:objective-mapping}
\begin{tabular}{|p{6cm}|l|p{5cm}|}
\hline
\textbf{Mục tiêu} & \textbf{Trạng thái} & \textbf{Ghi chú} \\
\hline
Làm rõ kiến thức nền tảng & Hoàn thành & Chương 2 \\
\hline
Đánh giá hệ thống Florus cũ & Hoàn thành & Mục 3.3 \\
\hline
Đề xuất giải pháp Kubernetes & Hoàn thành & Mục 3.4 \\
\hline
Thiết kế kiến trúc tổng thể & Hoàn thành & Chương 4 \\
\hline
Xây dựng Helm Charts & Hoàn thành & Mục 3.4.4 \\
\hline
Thử nghiệm trên Minikube & Hoàn thành & Mục 3.5 \\
\hline
% Huấn luyện mô hình AI & Chưa thực hiện & Chuyển sang ĐATN \\
% \hline
% Triển khai Production & Chưa thực hiện & Chuyển sang ĐATN \\
% \hline
\end{tabular}
\end{table}
